\par\noindent\textbf{(1)} 取两个50:50分束器的相位约定一致，输出单光子概率为 p₁=cos²(φ/2)、p₂=sin²(φ/2)，满足p₁+p₂=1。相干输入经线性光学仍为两个相干输出，故平均光子数 λ₁=Np₁、λ₂=Np₂。
\par\noindent\textbf{(2)} 相干态两输出计数 n₁,n₂ 为相互独立的Poisson变量，均值分别λ₁,λ₂，方差同值。计数差 D=n₁−n₂ 的均值 ⟨D⟩=Ncosφ，方差 Var(D)=λ₁+λ₂=N；其噪声与φ无关但差分均值的斜率为−Nsinφ。
\par\noindent\textbf{(3)} 两个独立Poisson通道的经典Fisher信息 F=Σ\_j(∂φλ\_j)²/λ\_j。代入 λ₁,λ₂ 得 F=N[sin²φ/(4cos²(φ/2))+sin²φ/(4sin²(φ/2))]=N（在非退化点按连续极限）。因此Cramér–Rao界 Δφ≥1/√N，为散粒噪声标度。
\par\noindent\textbf{(4)} 若只用差分均值线性误差传播，则 Δφ≈√N/(N|sinφ|)，在φ≈π/2最优；在暗端口角度直接线化会出现表观发散，但完整两端口似然仍含非高斯信息。单次强度只依赖cosφ，尚有全局φ与−φ模糊，需要已知偏置或多相位测量消歧。
